\documentclass[11pt,a4paper]{article}

\usepackage[utf8]{inputenc}
\usepackage[T1]{fontenc}
\usepackage{lmodern}
\usepackage[french]{babel}
\usepackage{amsmath}
\usepackage{booktabs}
\usepackage{listings}
\usepackage{xcolor}
\usepackage[hidelinks]{hyperref}

\lstset{
  basicstyle=\ttfamily\small,
  keywordstyle=\bfseries,
  commentstyle=\itshape\color{gray},
  showstringspaces=false,
  breaklines=true,
  frame=single,
  framesep=4pt
}

\begin{document}

\section{Théorie}

\section{Architecture}

\section{Implémentation FPGA}

\subsection{Bases}

\subsubsection{Quantification}

La quantification est nécessaire dans un réseau de neurones convolutif car
l'entraînement se fait à partir de valeurs représentées en float (32 ou
64~bits). Une multiplication de nombres (u)int8 peut être stockée en int32,
alors qu'une multiplication de float exige un matériel beaucoup plus coûteux
(alignement d'exposants, normalisation, etc.) (Jacob et al.). La
représentation en int8 est donc beaucoup plus adaptée aux « edge devices » :
sur le XC7A35T, un MAC int8$\times$int8 tient dans une fraction d'un DSP48.

La technique présentée dans Jacob et al. convertit un nombre réel $r$ à un
entier $q$ par une transformation affine :
\begin{equation*}
  r \approx S \cdot (q - Z)
\end{equation*}
où $S$ et $Z$ sont les paramètres de quantification, fixés par tenseur au
moment de la conversion :
\begin{itemize}
  \item $S$ (\emph{scale}) : nombre réel positif ;
  \item $Z$ (\emph{zero-point}) : entier du même type que $q$. C'est la
    valeur de $q$ qui correspond exactement à $r = 0$.
\end{itemize}

Exemple : représenter la plage réelle $[-1.0,\, 3.0]$ sur int8
($[-128,\, 127]$, 256~pas). L'échelle est
$S = (3.0 - (-1.0))/255 \approx 0.0157$, et le zero-point est choisi pour que
$q = -128$ corresponde à $r = -1.0$, soit $Z = -128 + 1.0/S \approx -64$.
On vérifie : $r = S \cdot (-64 - (-64)) = 0$.

Déquantification pour retrouver $r$ avec l'interpréteur TFLite :
\begin{lstlisting}[language=Python]
interp = tf.lite.Interpreter(model_content=tflite_model)
interp.allocate_tensors()
S, Z = interp.get_input_details()[0]["quantization"]
r = S * (q.astype(np.int32) - Z)                        # dequantize
\end{lstlisting}

\paragraph{Contraintes de la quantification int8 de TFLite}\mbox{}\\

La spécification 8~bits de TFLite/LiteRT a des particularités selon le rôle
du tenseur.

Tenseurs :
\begin{itemize}
  \item Activations : int8, quantification avec $Z$ variable (libre), un
    couple $(S, Z)$ par tenseur.
  \item Poids : int8, quantification avec $Z = 0$, un $S$ par canal de sortie
    (\emph{output channel}).
  \item Biais : int32, $S = S_{in} \cdot S_w$, $Z = 0$.
\end{itemize}

L'accumulateur int32 d'une couche est donc dans les unités
$S_{in} \cdot S_w$, alors que la couche suivante attend de l'int8 dans ses
propres unités $S_{out}$. Le passage de l'un à l'autre (multiplication par le
facteur effectif $M$, arrondi, ajout de $Z_{out}$, clamp) est la
\textbf{requantification}.

Le \textbf{facteur effectif} $M$ sort directement de l'arithmétique des
scales. En réel, un canal de sortie (\emph{output channel}) calcule :
\begin{equation*}
  r_{out} = \sum r_{in} \cdot r_w + r_{bias}
\end{equation*}
En remplaçant chaque $r$ par sa forme quantifiée $S \cdot (q - Z)$ (les poids
et le biais ayant $Z = 0$) :
\begin{align*}
  S_{out} \cdot (q_{out} - Z_{out})
    &= \sum S_{in} (q_{in} - Z_{in}) \cdot S_w\, q_w
       + S_{in} S_w\, q_{bias} \\
    &= S_{in} S_w \left[ \sum (q_{in} - Z_{in})\, q_w + q_{bias} \right] \\
    &= S_{in} S_w \cdot acc
\end{align*}
où $acc$ est exactement l'accumulateur entier du MAC. En isolant $q_{out}$ :
\begin{equation*}
  q_{out} = M \cdot acc + Z_{out},
  \qquad \text{avec} \qquad
  M = \frac{S_{in} \cdot S_w}{S_{out}}
\end{equation*}

$M$ est donc une constante qui condense les trois scales d'une couche. C'est
le « multiplicateur » défini par Jacob et al. éq.~(5),
$M := S_1 S_2 / S_3$, avec $S_1 = S_{in}$, $S_2 = S_w$, $S_3 = S_{out}$ dans
nos variables. Elle effectue la conversion entre les unités de
l'accumulateur et celles du tenseur de sortie. Il y en a un par canal de
sortie quand les poids sont quantifiés par canal. Comme $M$ est réel mais que
le FPGA ne calcule qu'en entier, il est appliqué via la décomposition
$M = m \cdot 2^{-31} \cdot 2^{-n}$ ($m$ entier Q31, $n$ = shift) : c'est la
forme normalisée $M = 2^{-n} \cdot M_0$ avec $M_0 \in [0.5,\, 1)$ de
l'équation~(6) du papier, où $m$ est l'entier int32 le plus proche de
$2^{31} \cdot M_0$. Ce calcul est détaillé aux sections « Export des
paramètres » et « requantize.v ».

Pour notre modèle : l'entrée a comme paramètres $S = 1/255$ et $Z = -128$, de
sorte que l'int8 d'entrée est simplement l'octet de la caméra XOR
\texttt{0x80}. La chaîne des zero-points d'activation actuelle est
$-128 \to -41 \to -81 \to -114 \to 62$ (conv1 $\to$ conv2 $\to$ dense1 $\to$
dense2).

\paragraph{TensorFlow}\mbox{}\\

Framework de construction, d'entraînement et d'export de modèles d'IA,
organisé autour du calcul sur tenseurs. Un tenseur est un tableau
n-dimensionnel typé :
\begin{itemize}
  \item rank-0 : scalaire
  \item rank-1 : vecteur (1D)
  \item rank-2 : matrice (2D)
  \item rank-3 et $+$ : tableaux 3D et plus
\end{itemize}
(\url{https://www.tensorflow.org/guide/tensor})

4 composantes de l'écosystème sont utilisées :
\begin{itemize}
  \item Keras : définition de l'architecture (Conv2D, MaxPooling2D, Dense) et
    entraînement float32 dans \texttt{training\_script.py} ;
  \item TensorFlow Model Optimization (tfmot) : le \emph{quantization-aware
    training} ;
  \item TFLiteConverter : convertit le modèle QAT en un fichier
    \texttt{.tflite} entièrement int8 : poids int8 par canal, biais int32, et
    les $(S, Z)$ de chaque tenseur ;
  \item LiteRT interpreter : extraction des paramètres des tenseurs en
    utilisant le nouvel API de Google.
\end{itemize}

\subsection{Training}

Nous utilisons la même architecture et outils qui ont été définis dans le
cahier des charges.

\subsection{Export des paramètres}

Le script \texttt{scripts/export\_params.py} fait la conversion des
paramètres entre le modèle entraîné et le RTL : il extrait du fichier
\texttt{asl\_cnn\_qaware.tflite} les poids, les biais et les paramètres de
quantification, puis génère les fichiers \texttt{.mem} chargés dans les ROM
du FPGA (via \texttt{\$readmemh}).

\subsubsection{Extraction depuis le .tflite}

L'extraction utilise l'interpréteur LiteRT
(\texttt{ai\_edge\_litert.Interpreter}, même API que
\texttt{tf.lite.Interpreter}). Les tenseurs sont repérés par leur index dans
le graphe. Ces indices ont été identifiés avec Netron.app et sont écrits dans
la table \texttt{LAYERS}, un quadruplet (entrée, poids, biais, sortie) par
couche :

\begin{center}
\begin{tabular}{lcccc}
  \toprule
  Couche & entrée & poids & biais & sortie \\
  \midrule
  conv1  & 5  & 6  & 2 & 7  \\
  conv2  & 8  & 9  & 1 & 10 \\
  dense1 & 12 & 13 & 4 & 14 \\
  dense2 & 14 & 15 & 3 & 16 \\
  \bottomrule
\end{tabular}
\end{center}

Pour chaque couche, \texttt{extract()} récupère sept valeurs. Elles
reprennent les paramètres $S$ et $Z$ du schéma $r \approx S \cdot (q - Z)$ de
la section « Quantification », appliqués aux trois tenseurs de la couche
(entrée, poids, sortie) :
\begin{itemize}
  \item \texttt{s\_in} : scale du tenseur d'entrée. Entre dans le calcul du
    facteur effectif $M = S_{in} \cdot S_w / S_{out}$.
  \item \texttt{zp\_in} : zero-point du tenseur d'entrée. C'est la valeur
    int8 qui représente le réel 0 en entrée ; le MAC la soustrait de chaque
    activation (\texttt{act - zp\_in} dans \texttt{mac\_unit.v}) pour
    ramener le calcul autour du vrai zéro.
  \item \texttt{s\_w} : scale(s) des poids. Tableau d'un seul élément si la
    couche a un scale global, ou d'un scale par canal de sortie
    (quantification par canal, cf. « Contraintes du schéma int8 de
    TFLite »). Les poids étant symétriques, ils n'ont pas de zero-point
    ($Z = 0$).
  \item \texttt{weights} : les poids int8 eux-mêmes, dans le layout du
    tenseur TFLite (conv : [canal de sortie, hauteur, largeur, canal
    d'entrée] ; dense : [sortie, entrée]). Aplatis tels quels dans
    \texttt{weights.mem}.
  \item \texttt{bias} : les biais int32, un par canal de sortie. Déjà
    exprimés dans les unités de l'accumulateur ($S$ forcé à
    $S_{in} \cdot S_w$, cf. contraintes TFLite), ils s'additionnent donc
    directement au premier terme du MAC sans conversion.
  \item \texttt{s\_out} : scale du tenseur de sortie. Dénominateur du facteur
    effectif $M$ ; c'est lui qui ramène l'accumulateur int32 dans l'échelle
    int8 de la couche suivante.
  \item \texttt{zp\_out} : zero-point du tenseur de sortie. Ajouté après la
    requantification (\texttt{relu\_clamp}) pour replacer le résultat dans la
    représentation du tenseur de sortie ; sert aussi de borne basse du clamp
    pour les couches ReLU.
\end{itemize}

Les scales ne sont pas dans les fichiers \texttt{.mem}. Il n'y a que les
poids, les biais et le couple $(m, \text{shift})$ qui encode $M$. Les
zero-points, eux, deviennent des constantes de \texttt{cnn\_params.vh}.

\subsubsection{Génération des fichiers .mem}

\texttt{convert\_mem()} écrit quatre fichiers dans
\texttt{vivado/conv\_accel/mem/}, toutes couches concaténées dans l'ordre
conv1 $\to$ conv2 $\to$ dense1 $\to$ dense2 :
\begin{itemize}
  \item \texttt{weights.mem} : poids int8 aplatis, un octet hexadécimal par
    ligne ;
  \item \texttt{bias.mem} : biais int32, 8 chiffres hexadécimaux par ligne ;
  \item \texttt{mult.mem} : multiplicateur de requantification $m$ (entier
    Q31, int32) ;
  \item \texttt{shift.mem} : décalage $n$ (int8), un couple $(m, n)$ par
    canal de sortie.
\end{itemize}

Le couple $(m, n)$ encode le facteur effectif
$M = S_{in} \cdot S_w / S_{out}$ de chaque canal (le $S_w$ par canal est
utilisé quand la couche en a plusieurs). C'est la décomposition
$M = m \cdot 2^{-31} \cdot 2^{-n}$ introduite à la section « Quantification »
et utilisée par \texttt{requantize.v}.

\subsubsection{quantize\_multiplier : conversion de float à Q31 + shift}

\paragraph{Le format Q31 (1.31)}\mbox{}\\

Q31 est un format à virgule fixe fractionnaire sur 32~bits, aussi noté 1.31 :
1~bit pour la partie entière, 31~bits pour la partie fractionnaire (spec
Microchip). Le bit de poids fort est le bit de signe (complément à deux) et
la virgule est implicite juste après lui. Un entier $q$ interprété en Q31
représente donc le réel :
\begin{equation*}
  r = \frac{q}{2^{31}}
\end{equation*}
\begin{itemize}
  \item Plage : $-1.0$ (\texttt{0x80000000}) à
    $1 - 2^{-31} \approx +0.9999999995$ (\texttt{0x7FFFFFFF}).
  \item Résolution : $2^{-31} \approx 4.66 \times 10^{-10}$ (le pas entre
    deux valeurs consécutives).
  \item Exemple (spec Microchip) : l'entier 1610612736
    (\texttt{0x60000000}) vaut 0.75 en Q31 (bits $2^{-1} + 2^{-2} = 0.5 +
    0.25$).
\end{itemize}

Dans notre cas, \texttt{mult.mem} stocke la mantisse normalisée
$M_0 \in [0.5,\, 1)$ en Q31 : $m = \operatorname{round}(2^{31} \cdot M_0)$,
donc $m \in [2^{30},\, 2^{31})$. $m$ est toujours positif (bit de signe à 0)
et son bit~30 est toujours à 1. Le format seul ne couvre que $[-1,\, 1)$ :
les valeurs de $M$ plus petites (jusqu'à $\sim 2^{-12}$) sont atteintes par
le décalage à droite additionnel $n$ (\texttt{shift.mem}), d'où le couple
« Q31 + shift ».

\paragraph{Justification de représentation}\mbox{}\\

$M$ est le seul nombre réel qui survit à la quantification en modèle TFLite :
tout le reste du calcul est déjà en entiers. Nous choisissons de le
représenter en Q31 + shift plutôt qu'en float pour répondre à
4~contraintes :
\begin{itemize}
  \item Coût matériel : le XC7A35T n'a pas de FPU. Une multiplication
    flottante serait beaucoup trop complexe. Avec Q31 + shift, la
    requantification se réduit à une multiplication entière
    $32 \times 32 \to 64$~bits suivie d'un décalage arithmétique, ce qui est
    beaucoup plus efficace en espace et en temps.
  \item Précision : un format fixed-point à virgule fixe unique gaspillerait
    des bits, car les $M$ du modèle s'étalent sur $\sim 2^{-7}$ à $2^{-12}$
    (d'où la plage du shift). La normalisation par \texttt{frexp} garantit
    une mantisse dans $[2^{30},\, 2^{31})$ quel que soit l'ordre de grandeur
    de $M$ : les 31~bits de précision sont toujours pleinement utilisés
    (erreur relative $\sim 2^{-31}$, négligeable devant le pas de
    quantification int8). C'est l'équivalent entier d'un flottant évalué sans
    hardware pour unités à points flottants.
  \item Espace de stockage : 32~bits ($m$) + un petit entier ($n$, 5~bits
    suffisent) par canal de sortie, soit 152 entrées dans
    \texttt{mult.mem}/\texttt{shift.mem}.
  \item Conformité : c'est exactement le schéma de TFLite/gemmlowp (Jacob et
    al., Parra et al.). Garder la même décomposition permet la comparaison
    avec la référence afin d'assurer une exactitude des opérations du FPGA.
\end{itemize}

\paragraph{Calcul}\mbox{}\\

TensorFlow calcule cette décomposition à l'intérieur de l'interpréteur
(\texttt{QuantizeMultiplier},
\texttt{tensorflow/lite/kernels/internal/quantization\_util.cc}) mais ne
l'expose pas comme appel de librairie ; la fonction est donc réécrite dans le
script :
\begin{enumerate}
  \item \texttt{math.frexp(M)} sépare $M$ en une mantisse $q \in [0.5,\, 1)$
    et un exposant de 2 ;
  \item la mantisse est mise à l'échelle Q31 ($q \cdot 2^{31}$) et arrondie
    \emph{half away from zero} (\texttt{round\_half\_away}, pour reproduire
    \texttt{TfLiteRound} et éviter le biais du round-half-even de Python sur
    un 0.5 exact) ;
  \item si l'arrondi pousse la mantisse à exactement $2^{31}$, elle est
    divisée par 2 et l'exposant incrémenté ;
  \item le shift retourné est $n = -\text{exposant}$ (convention décalage à
    droite, puisque $M < 1$).
\end{enumerate}

Ce calcul correspond exactement aux modules Verilog \texttt{srdhm}
(SaturatingRoundingDoublingHighMul) et \texttt{round\_shift} côté FPGA. Des
assertions vérifient les invariants supposés par le RTL : $M \in (0,\, 1)$,
multiplicateur normalisé ($2^{30} \le m < 2^{31}$) et shift dans la plage
$7 \le n \le 12$ vérifiée du RTL : un retrain qui sortirait de cette plage
ferait échouer l'export plutôt que de produire une inférence silencieusement
fausse.

\subsubsection{Constantes pour cnn\_params.vh}

Enfin, \texttt{write\_vh.py} (\texttt{print\_vh}) affiche les
\texttt{localparam} susceptibles de changer entre deux réentraînements, à
reporter manuellement dans \texttt{cnn\_params.vh} : pour chaque couche,
\texttt{ZP\_IN}, \texttt{ZP\_OUT} et \texttt{CLAMP\_MIN}. Le clamp bas vaut
$Z_{out}$ pour les couches suivies d'un ReLU fusionné (le zéro réel
correspond à $q = Z_{out}$), et $-128$ pour la dernière couche (logits bruts,
plage int8 complète).

\subsection{Modules}

\subsubsection{requantize.v}

Chaîne combinatoire en trois étages, représentative de l'Appendix~A de Parra
et al. :
\begin{enumerate}
  \item \texttt{srdhm} (\emph{SaturatingRoundingDoublingHighMul},
    Algorithm~A1) : produit 64~bits $acc \times m$, nudge d'arrondi ($+2^{30}$
    si positif, $1 - 2^{30}$ sinon), division par $2^{31}$ tronquée vers
    zéro. La saturation du cas $a = b = \text{INT32\_MIN}$ de l'algorithme
    est omise car $m$ est toujours positif.
  \item \texttt{round\_shift} (\emph{RoundingDivideByPOT}, Algorithm~A2) :
    arrondi de $q / 2^n$ avec la convention exacte de gemmlowp
    ($\text{threshold} = (\text{mask} \gg 1) + (q < 0)$) : les ties
    s'arrondissent en s'éloignant de zéro.
  \item \texttt{relu\_clamp} : ajout de $Z_{out}$ puis clamp à
    $[\text{CLAMP\_MIN},\, 127]$ (CLAMP\_MIN $= Z_{out}$ pour les couches
    ReLU, $-128$ pour les logits).
\end{enumerate}

\paragraph{Choix : double rounding (papier) vs single rounding (LiteRT
moderne)}\mbox{}\\

Le pipeline srdhm vers round\_shift est le \emph{double arrondi} historique
de gemmlowp (librairie de multiplication de matrices de Google), celui décrit
par Jacob et al. et Parra et al. TFLite/LiteRT utilisent maintenant, par
défaut, un \emph{single rounding} afin de régler quelques erreurs (voir
\url{https://github.com/tensorflow/tensorflow/issues/25087})
(\texttt{TFLITE\_SINGLE\_ROUNDING}) fourni par la librairie ruy, successeur à
gemmlowp. La formule résultante est
$r = (acc \cdot m + 2^{30+n}) \gg (31+n)$.

\subsubsection{mac\_unit.v}

Ce module multiplie et accumule le produit des activations avec les poids du
filtre fournis par l'entraînement. Multiplication : le produit est le
résultat de l'activation soustrait au zero-point, ensuite multiplié par le
poids. Addition : dans le cas du premier terme, le biais est additionné au
produit. Par après, le produit est additionné au total.

\subsubsection{conv\_engine.v}

Ce module est utilisé pour les calculs des couches conv1, conv2, dense1,
dense2. Il effectue la convolution, soit l'opération MAC ainsi que la
requantification.

Par exemple, pour 2 pixels de sortie :

Si \texttt{s\_valid} (alors \texttt{s\_valid} = 1), l'unité de MAC effectue
les opérations.

\subsubsection{pool\_engine.v}

Effectue MaxPooling sur les données.

\subsubsection{Mémoire}

\paragraph{param\_rom.v}\mbox{}\\

\paragraph{act\_ram.v}\mbox{}\\

\subsubsection{Assemblage des couches}

\section*{Sources}
\addcontentsline{toc}{section}{Sources}

\begin{itemize}
  \item Jacob et al., \emph{Quantization and Training of Neural Networks for
    Efficient Integer-Arithmetic-Only Inference}, CVPR 2018 : the scheme,
    QAT, fused ReLU. \url{https://arxiv.org/abs/1712.05877}
  \item LiteRT 8-bit quantization.
    \url{https://developers.google.com/edge/litert/models/quantization_spec}
  \item gemmlowp \texttt{fixedpoint.h} : reference
    SaturatingRoundingDoublingHighMul / RoundingDivideByPOT.
    \url{https://github.com/google/gemmlowp/blob/master/fixedpoint/fixedpoint.h}
  \item ruy \texttt{apply\_multiplier.cc}.
    \url{https://github.com/google/ruy/blob/master/ruy/apply_multiplier.cc}
  \item TFLite \texttt{common.h} / \texttt{quantization\_util.cc} :
    MultiplyByQuantizedMultiplier and QuantizeMultiplier.
    \url{https://github.com/tensorflow/tensorflow/blob/master/tensorflow/lite/kernels/internal/common.h}
  \item Parra et al., \emph{A Methodology and Open-Source Tools to Implement
    CNNs Quantized with TensorFlow Lite on FPGAs}, Electronics 2023 : the
    FPGA methodology, Appendix~A. \url{https://doi.org/10.3390/electronics12204367}
  \item Microchip, \emph{MPLAB XC-DSC Libraries Reference Manual},
    \S5.1.2.3 « Q31 (1.31) Format » : définition du format Q31.
    \url{https://onlinedocs.microchip.com/oxy/GUID-70ACD6B0-A33F-4653-B192-8465EAD1FD98-en-US-11/GUID-889092D2-E1E4-4491-ABDB-47CEEBA41E1D.html}
  \item \url{https://poloclub.github.io/cnn-explainer/}
\end{itemize}

\end{document}
